% =====================================================================
%  ch21-no-go-and-spec.tex
%  Structural no-go, and the specification of the missing machine.
%  Labels used: \LT (theorem) \LC (conditional) \LN (no-go) \LM (exact
%  missing step) \LE (numerical evidence). No statement in this chapter is
%  marked closed, and nothing here bears on Clay.
% =====================================================================
\let\codeprev\code
\renewcommand{\code}[1]{\texttt{\footnotesize\detokenize{#1}}}

\section{Structural no-go, and the specification of the machine that does not exist}
\label{ch:nogo-spec}

The two preceding chapters are the scalar face of the wall: one identity, four exact models, and a
certificate on a single plaquette. This chapter is the structural face, and its content is negative
in a precise sense. It is not that the campaign has failed to find the machine; it is that the
\emph{shape} of every candidate on record is excluded, and that the shape which survives can be
written down as a specification with four inputs, one of which is missing.

Three things are established below. (i) Two exact benchmark models separate the two coordinates the
campaign used as targets: $G''>0$ with $m=0$, and $m>0$ with $G''=0$. Neither coordinate implies
the other, and the necessity of $f_1^2$ fails on the same model (\S\ref{sec:ng-bench}). (ii) Every
prepared structure --- reflection positivity, K\"all\'en--Lehmann positivity, Osterwalder--Schrader
reconstruction, topology, block-spin renormalisation, large $N$ --- is a \emph{stage} and not a
\emph{machine}: each guarantees that the theory exists, that the measure is non-negative, that the
transfer operator is self-adjoint and that its spectrum lies in $[0,1]$, and none produces a bound
on $\lambda_2$ (\S\ref{sec:ng-stage}). (iii) One elementary average lemma explains why twenty
candidate roads died at the same filter, and that filter is itself the gap test
(\S\ref{sec:ng-a1}). What remains is a specification (\S\ref{sec:ng-spec}) and exactly one
unclosed sufficient condition (\S\ref{sec:ng-open}).

\textbf{Boundary, stated before anything is read as progress.} Nothing in this chapter proves a gap,
closes a Clay step, or shrinks the open cell. It records that the search for a handle cannot end in
a number, and it converts that record into the only form in which it is useful. \textbf{Clay
remains $0\%$.}

\subsection{Two benchmark counterexamples, stated before anything else}
\label{sec:ng-bench}

A candidate tool is a quantity attached to a model. The first question about it is therefore not
whether it is computable but whether the two quantities the campaign used as targets are functions
of one another. They are not, and two exactly solvable models decide this. Neither needs four
dimensions, lattice data, or a coupling window.

\paragraph{Benchmark (a): the harmonic oscillator $\otimes$ free mode ($M_6$) \LT.}
Take $H=\tfrac12(p_a^2+p_b^2)+\tfrac12\omega^2a^2$ on $L^2(\mathbb R^2)$: the $a$-direction is a
harmonic oscillator, the $b$-direction is free. Let $\rho=|\psi_0|^2$ be the marginal density of
$a$ and $V$ the potential in that direction. Using $-\tfrac12\psi_0''+V\psi_0=E_0\psi_0$,
\begin{equation}
G''(a)=-\partial_a^2\log\rho(a)
      =-2\,\frac{\psi_0''(a)}{\psi_0(a)}+2\Big(\frac{\psi_0'(a)}{\psi_0(a)}\Big)^{2}
      =4\big(E_0-V(a)\big)+2\Big(\frac{\psi_0'(a)}{\psi_0(a)}\Big)^{2}.
\label{eq:ng-m6}
\end{equation}
At the symmetric point $a=0$ the wave function is even, so $\psi_0'(0)=0$ and the second term
vanishes: $G''(0)=4(E_0-V(0))=2\omega$. The spectrum, however, is $\mathrm{spec}(H)=[\omega/2,\infty)$
and continuous, because the free direction contributes $[0,\infty)$ and the oscillator contributes
$\{\omega(n+\tfrac12)\}$; hence $m=0$. With $\omega=1$: $G''(0)=2.000000>0$ and $m=0$, both exact.

Read \eqref{eq:ng-m6} literally: $G''(0)>0$ is equivalent to $V(0)<E_0$, that is, to the origin
lying above the ground-state energy --- in the classically forbidden region, at a barrier top, or
equivalently at a local maximum of the marginal density. It measures \emph{whether the origin lies in
the allowed region}. It is a statement about one point. The gap $m=E_1-E_0$ is a statement about the
whole spectrum, set by the global shape of the potential and by the boundary conditions. There is no
function relating them, and the same identity reproduces the double-well value
$4(E_0-V(0))=0.5889397$ against the measured $0.5889$ (difference $3.9\times10^{-7}$, grid error).
\src{\code{T03-TOY.md} \S2 ($M_6$), \S3 (the identity and its two cross-checks: $2.000006$ for the
oscillator, $3.9\times10^{-7}$ for the double well); \code{U09-M6.md} \S1, \S5 for the reading
$m_L=\min(\omega,\tfrac12(2\pi/L)^2)\to0$; reproduction \code{research/v3/w7/t03_toy_models.py}.
\LT{} at the level of the identity and of the spectrum; the model is $0+1$-dimensional quantum
mechanics, not a gauge theory.}

\paragraph{Benchmark (b): two-dimensional Yang--Mills \LC.}
On a spatial circle the transfer operator of 2D $\mathrm{SU}(2)$ Yang--Mills is diagonal in the
character basis, $T=\sum_j w_jP_j$ with $w_j=e^{-\lambda C_2(j)}$, $\lambda=g^2/2$ and
$C_2(j)=j(j+1)$. The ground state is the trivial representation, i.e.\ the constant function, so
\emph{every} coarse marginal is flat and $\partial^2\log\rho\equiv0$: $G''(0)=0$ exactly, at every
coupling. The spectrum is $E_j=\lambda C_2(j)$, so
\begin{equation}
m=\lambda\,C_2(\tfrac12)=\tfrac34\lambda=\tfrac38g^2=0.375\qquad(g^2=1),
\label{eq:ng-2dym}
\end{equation}
which is positive. The same model kills the necessity of $f_1^2$: in two dimensions
$\operatorname{tr}F^2=H$ along the transfer direction is conserved (there is no propagating degree of
freedom), hence $O|\Omega\rangle\propto E_0|\Omega\rangle$ and
$\langle\Omega|O|\mathbb G\rangle=E_0\langle\Omega|\mathbb G\rangle=0$; character orthogonality gives
the same conclusion. Thus $f_1^2\equiv0$ in a theory with a positive gap.
\src{\code{T03-TOY.md} \S2 ($M_2$), labelled there \code{[T]*}: exact inside the model, the
literature not read verbatim; the paper keeps the weaker label. $f_1^2\equiv0$:
\code{B05-INVENT.md} \S1, \code{D05-FKG.md} \S3.4, \code{IMP-ARROWS.md} \S\S3, 5.}

\begin{claimbox}{cN}
\claimhead{Corollary}{(two directions of separation)}\label{cor:ng-sep}
In the class of models where both quantities are defined, neither implication holds:
$G''(0)=2\omega>0$ with $m=0$ ($M_6$), and $m=3g^2/8>0$ with $G''(0)=0$ (2D Yang--Mills).
Hence $G''(0)>0$ is not sufficient for $m>0$, and $m>0$ does not force $G''(0)>0$. The necessity
of $f_1^2$ fails separately, on the second model: a positive gap with $f_1^2\equiv0$.
\end{claimbox}

\paragraph{What these counterexamples do not say.}
They do not refute the identity layer. At fixed $(L,a)$ within pure Yang--Mills,
$G''(0)=E[H+\Phi''\mid0]-\operatorname{Var}(G+\Phi'\mid0)=V_{\mathrm{eff}}''=m_{\mathrm{eff}}^2(k\to0)$
is exact, so a counterexample \emph{inside} pure Yang--Mills would have to make the infimum over
$(L,a)$ fail --- which is the wall, not a refutation of an implication. What fails is the bare
implication in the natural class of reflection-positive lattice gauge theories: for
$G_{\mathrm{YM}}\times U(1)_{\mathrm{weak}}$ the Yang--Mills factor is gapped while the photon sector
is gapless, so $G''>0$ and $m=0$ coexist there as well. The correct reading is therefore the one the
ledger already uses: $G''(0)$ is a \emph{criterion} (an equivalent form of $H1$ at fixed $(L,a)$), not
a \emph{signpost}.
\src{\code{T01-IMPLY.md} \S1 (four names of $G''$), \S4.1 (the three-step chain, two of whose steps
are renamings), \S4.3 (the class-level counterexample and why it does not relocate the route);
\code{docs/paper-v2/ch05-theorems.tex} \code{eq:split}; \code{ch08-immune.tex} for the four
equivalent names.}

\subsection{The four-model touchstone}
\label{sec:ng-touchstone}

From here on no candidate is discussed before it is scored on Table~\ref{tab:ng-bench}. The table has
four rows and, as a test, only two bits of information.

\footnotesize
\begin{center}
\begin{tabular}{@{}lcccc@{}}
\toprule
Model & $m$ & $G''$ & $f_1^2$ & any new tool must give\\
\midrule
free massless field & $0$ & $0$ & $0$ & $\mathbf 0$\\
\textbf{2D Yang--Mills} & $\mathbf{>0}$ & $0$ & $\equiv0$ & $\mathbf{>0}$\\
\textbf{$M_6$} & $\mathbf 0$ & $\mathbf{2\omega>0}$ & --- & $\mathbf 0$\\
single plaquette & undefined (no transfer) & $+0.136$ & --- & \textbf{undefined}\\
\bottomrule
\end{tabular}
\captionof{table}{The touchstone. The single-plaquette entry is the certificate value
$G''(0)=+0.135994711041988098$ of Chapter~\ref{ch:h1-status}, reproduced there against its anchor
to $1.1\times10^{-16}$; that model has no transfer direction, so its $m$ is \emph{undefined} and
neither $0$ nor positive. The last column is the requirement, not a measurement.}
\label{tab:ng-bench}
\end{center}
\normalsize

Two remarks fix the table's force, and both are negative.

First, the discipline in the fourth row is load-bearing. A zero-dimensional model has no
$E_1$; writing ``undefined'' as $0$ manufactures a counterexample, and writing it as ``$>0$''
defends $H1$ for free. The single-plaquette certificate is a strict positivity statement of the
left half of $H1$ and was never evidence for the right half.

Second, the table cannot certify that a passing candidate is \emph{not} a relabelling. Its rows can
report only (i) whether a transfer direction exists and (ii) whether the floor is zero; the signature
$(0,+,0,\text{undefined})$ is therefore shared by $m$ itself and by every monotone relabelling of $m$.
Three candidates that pass all four rows are known --- the string tension $\sigma$, the Kato isolation
margin, and the Hilbert--Schmidt decay rate --- and the third is provably $2\lambda_{\mathrm{bot}}$,
i.e.\ a renaming, while the second equals $m$ on any relativistic target. The passing set is a subset
of the monotone relabellings of $m$.
\src{\code{T03-TOY.md} \S1 (the undefined-value discipline); \code{B05-INVENT.md} \S0, \S4 (B05-1:
the test carries two bits; all four candidates collide with $m$), \S3 (the three candidates and their
death causes). \LT{} for the undefined-value discipline and for B05-1; the numerics are \LE.}

\subsection{The structural no-go: every prepared structure is a stage, not a machine}
\label{sec:ng-stage}

The campaign did not lack structures. It had six, all standard, all applicable, and each of them is
quoted below with the number that shows it is a stage.

\begin{enumerate}[label=(\roman*)]
\item \textbf{Reflection positivity} \LT. It gives a self-adjoint contraction $T=e^{-aH}$,
  $0\le T\le1$, $\mathrm{spec}\,T\subset[0,1]$ (the negative half-axis is removed), a unique vacuum
  $\lambda_1=1$ simple, and $T>0$ pointwise at finite volume. Its entire quantitative content about
  the gap is $\lambda_2\le1$, i.e.\ $m\ge0$. Worse, it holds on both benchmark models: on 2D
  Yang--Mills the diagonal weights satisfy $w_j=e^{-\lambda C_2(j)}\ge0$ for the whole family
  $\lambda\ge0$, so every reflection-positivity inequality holds pointwise while $m(\lambda)=3\lambda/4$
  runs from $0$ to $3.75$; on $M_6$ the two-point kernel is an explicit square,
  $\langle\Theta f,f\rangle=(2\omega)^{-1}(\int e^{-\omega t}f_a)^2\cdot(\text{free-direction square})\ge0$
  (checked to $1.8\times10^{-15}$), with $m=0$. A structure that holds on a family whose gap varies
  from $0$ to $3.75$ cannot determine the gap.
\item \textbf{K\"all\'en--Lehmann and Osterwalder--Schrader positivity} \LT. These guarantee
  $\nu\ge0$, hence that $C(t)=\int e^{-Et}d\nu$ is completely monotone and that
  $m=\inf\operatorname{supp}\nu$ is a rate rather than a complex pole. They are why the word
  ``spectrum'' can be used at all. They supply no inequality in $\lambda_2$; the four readings of
  $m$ that they make available --- $-\log\lambda_2$, $\inf\operatorname{supp}\nu$, twice the
  Poincar\'e constant, $-\frac1a\log\|T|_{\Omega^\perp}\|$ --- are four names of the target.
\item \textbf{Topology} \LN. In two dimensions $F\wedge F\equiv0$ \emph{identically}: $F\in\Omega^2$
  and $F\wedge F\in\Omega^4=0$ on a two-manifold, and $\pi_1(\mathrm{SU}(N))=0$ makes the spectral
  flow vanish as well. Every topological reading is therefore $0$ on the one exactly solvable gapped
  gauge theory, and a $\beta$-independent quantity that is constant on connected components can only
  fall in one of two branches: it vanishes on 2D Yang--Mills (where $m>0$), or it is non-zero on a
  single configuration (which has no transfer matrix and hence no $m$). There is no third branch.
\item \textbf{Block spin} \LN. One step is a weakening, not a strengthening:
  $\beta_{W,\mathrm{eff}}=\beta_W-C_N\ln b$ with $C_N=11N^2/12\pi^2>0$, giving
  $\beta_{\mathrm{eff}}/\beta=0.897$ at the $\mathrm{SU}(2)$ crossover $\beta_W=2.5$ and $0.898$ at
  the $\mathrm{SU}(3)$ crossover $5.69$. The reason no power-law improvement exists in four
  dimensions is an exponent, not a constant: $\beta_W=2N/(g^2a^{4-d})$ gives engineering exponent
  $4-d$, which is $0$ at $d=4$ --- exactly marginal, logarithmic flow, sign fixed by the one-loop
  coefficient.
\item \textbf{Large $N$} \LN. The four-dimensional $N=\infty$ gap is an exact missing step, not a
  theorem; the available proxy $\rho^2=\Lambda_\star/m_1^2=2$ is also the value of the free field,
  where $e_1=1$, $e_2=2$ on $\mathbb T^4$ give $\rho^2=2$ exactly. A discriminant that takes the
  same value on a gapped theory and on a free field is not a discriminant.
\end{enumerate}

The list can be summarised in one step, and this is the chapter's main negative statement.

\begin{qbox}
\textbf{No handle exists: every handle is either the gap itself or blind to it.} \LN, named class:
quantities drawn from the prepared structures --- reflection positivity, K\"all\'en--Lehmann/OS
positivity, topology (indices, winding, spectral flow), block-spin renormalisation, and the
$1/N$ expansion.
\end{qbox}

The reason is a dichotomy that the corpus reached from two independent directions, and it is worth
stating as arithmetic rather than as opinion. A \emph{lossless} quantity --- $\lambda_2$, the
Poincar\'e constant, $\inf\operatorname{supp}\nu$, $f_1^2$ --- is equivalent to the gap, so a machine
built on one of them is the gap renamed. A \emph{lossy but inconsistent} quantity --- the twenty
candidates of \S\ref{sec:ng-a1} --- reports either a positive value on $M_6$ (a false certificate) or
zero (empty). A machine must therefore be lossy \emph{and} consistently positive \emph{and} of bounded
loss: an output $\gamma$ with
\begin{equation}
\gamma\ \ge\ \rho(T|_{\Omega^\perp})=e^{-a\,m_{\mathrm{eff}}}
\qquad\text{and}\qquad
\gamma\ \le\ e^{-a\Delta}<1\ \ \text{uniformly}.
\label{eq:ng-lossy}
\end{equation}
The output is strictly worse than $\lambda_2$, and by an unknown factor. That is exactly why it can
be proved, and it is also the operational test for ``not a renaming''.
\src{\code{B01-RP.md} \S1 (the five outputs of reflection positivity, two of them quantitative),
\S3.1 (the 2D family and the $0\to3.75$ range), \S3.2 (the $M_6$ square, $1.8\times10^{-15}$),
\S3.4 (the blindness statement), \S5 (what is missing); \code{B04-INDEX.md} \S4 (the
$\beta$-independent dichotomy), \S5.2; \code{B03-BLOCK.md} \S2.1 ($\beta_{\mathrm{eff}}/\beta=0.897$),
\S2.3 ($4-d=0$); \code{B02-LARGEN.md} \S3 (the free-field $\rho^2=2$); \code{D06-SPEC.md} \S1.5
(the lossy/lossless dichotomy and \eqref{eq:ng-lossy}).}

\subsection{Lemma A1, and why twenty roads died at the same filter}
\label{sec:ng-a1}

Let $d\nu\ge0$ on $[0,\infty)$ with $0<\int d\nu<\infty$ and put
$m_1:=\inf\operatorname{supp}\nu$. For measurable $w\ge0$ with $0<\int w\,d\nu<\infty$, define the
$s$-weighted average $A_w(\nu):=\big(\int s\,w(s)\,d\nu(s)\big)/\big(\int w(s)\,d\nu(s)\big)$.

\begin{claimbox}{cT}
\claimhead{Lemma}{(A1, the average lemma)}\label{lem:ng-a1}
$A_w(\nu)\ge m_1(\nu)$, with equality if and only if $w\cdot\nu$ is a point mass.
\end{claimbox}

The argument is two lines: $s\ge m_1$ pointwise on $\operatorname{supp}\nu$, and a non-negative weight
preserves order. Two corollaries carry the campaign's accounting.

\begin{claimbox}{cT}
\claimhead{Corollary}{(A1$'$ and A1$''$)}

(A1$'$, direction) Every average-type functional --- finitely many polynomial moments
$\int s^k\,d\nu$, any finite-$t$ regularisation or truncation, any $s$-extremum such as a
Cheeger constant, a capacity or a large-deviation rate --- is an \emph{upper} bound:
$\mathfrak F\ge m_1$. No such functional can prove $m_1\ge\delta$.

(A1$''$, the $M_6$ killer) If $m_1=0$ while $\nu\neq\delta_0$ --- a weight without a floor, which is
exactly the position of $M_6$ --- then $A_w>0$. Every average-type functional is therefore positive
on $M_6$, so filter~1 of the touchstone ($M_6\mapsto0$) must fail.
\end{claimbox}

The consequence for the ledger is a single sentence with twenty instances. All twenty candidates of
the species waves compute an $s$-weighted average or an $s$-extremum; by A1$'$ they bound the floor
from above, by A1$''$ they are positive on $M_6$, and therefore \textbf{all twenty die at filter~1}
--- not at the output end, not for want of a better constant. And filter~1 is itself the gap test: a
functional that separates $0$ from $>0$ on the pair
$(M_6,\text{2D YM})$ is, verbatim, a floor detector. The criteria therefore presuppose the
conclusion.
\src{\code{D06-SPEC.md} \S1.1 (A1 and its two corollaries), \S1.2 (the mechanism is averaging, not
monotonicity; the real obstruction is weak continuity), \S1.3 (the twenty-road accounting), \S1.4
(the repair: read ``whole'' on the model family, not on spacetime). \LT{} for A1 and both corollaries.}

One repair is available and its price must be recorded with it. If ``whole'' is read as
\emph{whole on the spacetime}, then the functional is an average and A1 applies. If ``whole'' is read
as \emph{whole on the model family} --- the domain becomes $\{\text{family}\}\times\{\beta\text{-interval}\}$
instead of a single model --- then filter~1 stops being a pointwise predicate, a
certificate machine may honestly say ``cannot certify'' on a single model, and A1 no longer applies.
The price is that the machine no longer returns a number: it returns a modulus and a proof
(\S\ref{sec:ng-spec}).

\subsection{Five independent impossibilities on the transport side}
\label{sec:ng-transport}

The transport problem asks for $f_1^2(2.5)>0$ given the seed $f_1^2(\beta_1)\ge r_0>0$ at
$\beta_1<1/48$. Five refutations are on record; each is independent of the others, and each closes a
different route.

\begin{enumerate}[label=(\roman*)]
\item \textbf{The provable lower bound on the vacuum overlap is $e^{-\Theta(V)}$, and it is tight.}
  \LT{} Write $\Omega=|\langle\psi_\beta|\psi_{\beta'}\rangle|^2=BC(\mu_\beta,\mu_{\beta'})^2$, the
  squared Bhattacharyya/Hellinger affinity of the two Perron--Frobenius measures. Then
  $\Omega\ge1/(n^2R_\beta R_{\beta'})$ with $n=e^{\Theta(V)}$ and $R_\beta=e^{\Theta(V)}$, so the only
  provable bound is $e^{-\Theta(V)}$. It cannot be improved, because the free-field value is of the
  same order: $\Omega_k=2\sqrt{\beta\beta'}/(\beta+\beta')=0.1810653082661706$ at
  $(\beta,\beta')=(2.5,1/48)$, independent of the mass, and multiplicativity gives
  $\Omega=0.1811^{\,n_{\mathrm{dof}}}$, i.e.\ $\log_{10}\Omega\approx-36479$ for four-dimensional
  $\mathrm{SU}(2)$ at $L=8$ with $n_{\mathrm{dof}}=3\cdot4\cdot L^4=49152$.
\item \textbf{$\Omega=1$ is literally false.} \LC{} Hellinger is an $f$-divergence and therefore
  monotone under channels; taking the marginal on a single spatial plaquette $P_0$ gives
  $\Omega\le\Omega_{P_0}$, so $1-\Omega\ge1-\Omega_{P_0}>0$ independently of the volume. The
  inequality is conditional on the $\beta=2.5$ plaquette distribution, i.e.\ on a non-perturbative
  input, and it is labelled accordingly. The opposite failure is exact and points the same way: on
  2D Yang--Mills and on a single plaquette the Perron--Frobenius vector is $\beta$-independent
  ($\psi\equiv\mathrm{const}$), hence $\Omega=1$ \emph{exactly} while the atom is $0$. $\Omega$
  compares vacua; the transport compares residues.
\item \textbf{No sign is available on the residue side.} \LT{} Two independent statements.
  (a) \emph{Componentwise monotonicity of $|\mathbb G_\beta\rangle$ is not merely unproved, it is not a
  well-posed statement}: with the convention $\Omega\equiv1$, $\lambda_0=1$ one has
  $\langle\mathbb G_\beta|1\rangle=0$, so $\mathbb G_\beta$ changes sign and its nodal set is non-empty;
  if in some phase convention $\mathbb G_\beta\ge\mathbb G_{\beta'}$ pointwise, equal norms force
  $\mathbb G_\beta=\mathbb G_{\beta'}$, two different kernels would share an eigenvector, and
  $\Phi\equiv\mathrm{const}$, a contradiction; and the overall phase is a convention under which the
  assertion flips.
  (b) \emph{The connected operator is sign-indefinite at every $\beta$}: because
  $\langle 1|\mathbb G_\beta\rangle=0$, the functional $F=\langle1|O|\mathbb G_\beta\rangle$ is
  invariant under $O\mapsto O-c$ for every constant $c$; choosing $c=\langle O\rangle$ --- the
  subtraction physics requires --- makes the numerical range straddle zero. Positivity of $O$ as a
  plaquette operator is therefore unusable, and it is not a question of a missing inequality.
\item \textbf{Gr\"onwall's hypothesis is verbatim its conclusion.} \LT{} For $F\in C^1([\beta_1,2.5])$
  with $F(\beta_1)>0$, the following are equivalent: $\exists C<\infty$ with $|F'|\le C|F|$;
  $F>0$ on the interval; $\log F$ Lipschitz. The implication that looks like the free step --- first
  zero argument --- is what proves it; and at a zero $\beta^*$ the two-dimensional solution forces
  $f_1^2\approx c(\beta-\beta^*)^2$, so $|F'/F|=2/|\beta-\beta^*|+O(1)$ and $C=\infty$. The existence
  layer fails too: writing $d|\mathbb G\rangle/d\beta=\sum_{n\neq0}|n\rangle\langle n|S_W|\mathbb G\rangle/(E_0-E_n)$
  already presupposes $E_0$ isolated and simple on the whole interval, which is the target.
\item \textbf{The transfer operator at $L=2$ is infinite-dimensional.} \LT{} At $L=2$ the slice
  carries $3L^3=24$ spatial links, and $T$ acts on $L^2(\mathrm{SU}(2)^{24})$, a separable
  infinite-dimensional Hilbert space; the $143$ of the certificate ledger is not its dimension but the
  gauge-quotient count of the full four-dimensional torus, $4L^4\cdot3-3L^4-1=192-48-1$. The
  matrix-element integral is $24$-dimensional, and a Peter--Weyl truncation at $j\le J$ gives
  $b_J=\sum_{m\le 2J+1}m^2$ states per link ($J=3$: $140$), hence $140^{24}\approx10^{51.5}$ states at
  $L=2$ and $10^{412}$ at $L=4$. Direct certification at the physical point is therefore not a matter
  of more computing, and positivity of the kernel, while it gives $m(L)>0$ for each $L$, gives no
  number.
\end{enumerate}
\src{\code{D01-OMEGA.md} \S2--\S4 ((i) and (ii): the bound, the tightness, the two failure
directions); \code{D05-FKG.md} \S2 ((iii)(a): the nodal argument, the equal-norm argument, the phase
objection), \S3.2 ((iii)(b): the translation freedom); \code{D04-ODE.md} \S1.6, \S2.1, \S2.3 ((iv):
the equivalence, the divergence rate, the existence layer); \code{D02-DIRECT.md} \S1, \S3.1, \S4 ((v):
the dimension count, the state count, the per-$L$ versus uniform distinction; reproduction
\code{w7/d02_dimcount.py} and \code{w7/d02_cert.py}).}

\subsection{The specification of the missing machine}
\label{sec:ng-spec}

Everything above is an exclusion. What follows is the positive residue: the object that is missing,
written as a specification with acceptance criteria, so that it can be attacked and so that its
failure can be diagnosed.

\paragraph{Type, and the trap in the output convention.}
The machine is not a functional $\mathfrak F:\text{Model}\to\mathbb R$; on that type the no-go of
\S\ref{sec:ng-a1} applies. It is a map $\mathsf M:(\mathfrak M,\mathfrak P,[\beta_0,\beta_1])\to(\Delta,\text{proof})$
from a model family and a structural packet to a modulus and a proof, with the generator convention
fixed once and for all:
$\lambda_2^{\mathrm{gen}}:=m_{\mathrm{eff}}^2\ge0$, and the machine's output is
$\lambda_2^{\mathrm{gen}}\ge\Delta^2>0$, equivalently
$\|T|_{\Omega^\perp}\|=e^{-a\,m_{\mathrm{eff}}}\le e^{-a\Delta}<1$. The inequality $\|T|_{\Omega^\perp}\|\le1$
is free from reflection positivity; only \emph{strictness with an exponent} is content. Reading
$\lambda_2$ as $\lambda_2(T)$ reverses the required direction, and mixing the two conventions
manufactures progress.

\paragraph{Inputs.}
Four, of which three exist. (I1) the family $\{\mu_{\beta,L,a}\}$ with a $\beta$-interval containing
$2.5$, as a statement uniform on the family rather than per $(L,a)$: available from the
Kogut--Susskind/Wilson formulation. (I2) a compactness witness --- a uniform \emph{a priori} bound
making the family precompact in the Käll\'en--Lehmann weak topology (a uniform bound on
$\sup_p\langle\tfrac1N\operatorname{Re}\operatorname{tr}U_p\rangle$, or an entropy/BV bound): only
partially available, hence \LM. (I3) a seed,
a proved strict floor at the left endpoint: $f_1^2>0$ for $\beta_0<1/48$, unconditionally \LT. (I4)
the object itself: an explicit positive operator, weight or isomorphism carrying the certificate,
checkable by finite-dimensional truncation on $\ell$-blocks. (I4) is what is missing.

\paragraph{Three operational criteria for ``not a rename''.}
(1) \emph{Visible loss}: the output and the truth differ by a factor that cannot be removed; a
candidate claiming $\gamma=\rho(T)$ is a renaming. (2) \emph{The certificate object contains no
spectrum}: the carrier is algebraic or geometric, not a $\sup/\inf$ of a spectrum. (3) \emph{Failure is
diagnosable}: at a given capacity the candidate must emit a dual obstruction certificate of Farkas
type explaining why that capacity is insufficient. None of the twenty candidates of
\S\ref{sec:ng-a1} has property (3); on failure they report only ``the value is $0$''.

\paragraph{Five filters, in the machine reading.}
The touchstone of \S\ref{sec:ng-touchstone} becomes a scoring rule for a machine rather than for a
number, and the four rows become four behaviours rather than four values.

\footnotesize
\begin{center}
\begin{tabular}{@{}p{0.6cm}>{\raggedright\arraybackslash}p{5.5cm}>{\raggedright\arraybackslash}p{7.4cm}@{}}
\toprule
\# & Filter & Machine reading\\
\midrule
1 & four models & the machine \emph{refuses} to certify on the free field, on $M_6$ and on a single plaquette, and it \emph{must truly build} a certificate on 2D Yang--Mills --- if it cannot, it is empty\\
2 & not of $\min$ type & the certificate is an operator or supremum statement ($\|\cdot\|$, $\succeq$), not an $\inf$ over a set\\
3 & whole & whole \emph{on the model family}: one certificate covers a whole $\beta$-segment; this is the repaired reading of \S\ref{sec:ng-a1}\\
4 & not a constant & $\Delta=\Delta(\beta,\Lambda)$, rebuilt at each $\beta$\\
5 & dimensions & $\Delta$ carries mass dimension and $g\to0$ forces $\rho\to1$, $\Delta\to0$ (it must degenerate when $W\to1$)\\
\bottomrule
\end{tabular}
\captionof{table}{The five filters as a scoring rule for a machine. Filter~1 is a
non-emptiness test, not a positivity test: a candidate that cannot be built on a theory already known
to be gapped is not a machine.}
\label{tab:ng-filters}
\end{center}
\normalsize

\paragraph{Three candidate designs, one of them viable.}
\LT{} is not claimed for any of them; the entries are a classification of the design space.

\footnotesize
\begin{center}
\begin{tabular}{@{}p{2.5cm}>{\raggedright\arraybackslash}p{4.4cm}>{\raggedright\arraybackslash}p{2.4cm}>{\raggedright\arraybackslash}p{3.4cm}@{}}
\toprule
Candidate & Mechanism class & Filters & Exact missing object\\
\midrule
D1 \textbf{W-SHIFT} & similarity transform plus one operator inequality; contains the dead
Bakry--\'Emery route as the single point $c=\tfrac12$, $K=S$; not variational (the variational
direction is opposite) & 5/5 & \code{W-CALC}\\
D3 \textbf{OS-SCHR\"O} & Hilbert-space isomorphism to $L^2(\mathcal A/\mathcal G,d\mu_W)$, then
$T^2-mT\succeq0$; needs no $E_0$ because $T\succeq0$ by construction & 5/5 \emph{formally} &
\code{KKN4}\\
D2 \textbf{RETRACT} & topology of the parameter space (closed floor set, non-empty seed, path lift) &
5/5 & \code{RETRACT}$(\delta)$\\
\bottomrule
\end{tabular}
\captionof{table}{The three candidates. D1 is ranked first: it has a three-dimensional precedent, it
strictly contains a dead route as a special case rather than renaming it, and its failure mode is
diagnosable. D2 is logically legal but the missing object \code{RETRACT}$(\delta)$ is equivalent to a
gap on the interval, so it relocates the wall to a single point rather than removing it. D3 has the
hardest object: no four-dimensional orbit-space representation is known, while the three-dimensional
case exists (Nair 2026, \code{arXiv:2608.10133}, abstract level, not read verbatim).}
\label{tab:ng-cands}
\end{center}
\normalsize

The specification of D1, verbatim, is this: for each $\beta$, find $(c,K)$ with
\begin{equation}
W^{1/2}\big(H-\langle\Omega,H\Omega\rangle\big)W^{-1/2}\ \succeq\ \Delta^2\,\mathbf 1_{\Omega^\perp},
\qquad W=e^{cK},
\label{eq:ng-d1}
\end{equation}
where $H$ is the Kogut--Susskind Hamiltonian and $K$ runs over a class of explicit, gauge-invariant,
local and explicitly squarable operators: the magnetic Hamiltonian
$K=\sum_p(1-\tfrac1N\operatorname{Re}\operatorname{tr}U_p)$, the plaquette sum
$K=\sum_p\tfrac1N\operatorname{Re}\operatorname{tr}U_p$, or a one-block local polynomial. The test is a
small-matrix computation in the electric-Casimir-truncated strong-coupling basis, and it must pass two
things: it must \emph{build} the certificate on 2D Yang--Mills (otherwise the machine is empty), and it
must \emph{fail} on the free field and on $M_6$ (otherwise the certificate is false). The missing object
is \code{W-CALC}, a weight calculus that cancels the Savvidy negative modes under the conjugation by
$W^{1/2}\cdot W^{-1/2}$ \emph{pointwise} rather than in integral form, and it must satisfy five
properties: $W=e^{cK}>0$ local and gauge-invariant; the negative modes of $W^{1/2}HW^{-1/2}$ cancelled
pointwise in the Savvidy directions; $\Delta\ge c_0>0$ independent of $(L,a)$; $\Delta\to0$ as $g\to0$;
and a Farkas-type obstruction on failure.
\src{\code{D06-SPEC.md} \S2.1--\S2.5 (type, four inputs, output convention, three criteria, five
filters), \S3.1 (D1 and its five filters), \S3.2 (D2), \S3.3 (D3), \S3.4 (four candidates killed on
the spot: PAD\'E-NEV and LOW-WEIGHT-COUNT by A1, ANOM-WARD because the anomaly is a contact term,
CERT-RG by the marginality $\beta_{\mathrm{eff}}/\beta=0.897$ and $4-d=0$), \S4.1 (the specification
\eqref{eq:ng-d1}), \S4.3 (the five properties of \code{W-CALC}). \LM{} for I2, I4 and all three
missing objects.}

\paragraph{The discipline, in three steps, and the theorem that replaces ``no value''.}
(1) Build the certificate on 2D Yang--Mills first; if it cannot be built there, the candidate is empty
and the verdict is not ``needs more computing''. (2) Quantify the wall as a curve:
$\ell^*(\beta):=\min\{\ell:\ \text{an } \ell\text{-block local }(c,K)\text{ closes the certificate}\}$,
computed at $L=2$, $\mathrm{SU}(2)$, with $j\le\tfrac12$ or $1$ --- four points
$\ell^*(0.5),\ell^*(1),\ell^*(2),\ell^*(2.5)$; if $\ell^*$ diverges with $\beta$, the wall is a
statement about the non-locality a certificate requires, which is a publishable diagnostic and not a
gap. (3) On failure, solve the dual (Farkas) problem and report \emph{which} direction --- transverse
magnetic Savvidy, zero mode, boundary --- blocked $(c,K)$; this is what upgrades ``the value is $0$''
into ``this dimension of the weight is missing''. And where no candidate can exist, the accounting is
itself a theorem:

\begin{claimbox}{cN}
\claimhead{Theorem}{(D06-T3, filter incompatibility)}\label{thm:ng-d06t3}
If $\mathfrak F$ is built from an $s$-weighted average or an $s$-extremum of $\nu$ --- which covers all
polynomial moments, all finite-$t$ regularisations, and all capacity- or rate-type extrema --- then
$\mathfrak F(\nu)\ge m_1(\nu)$ by A1, and $\mathfrak F(\mu_{M_6})>0$ on every model with $m_1=0$ and
$\nu\neq\delta_0$. Hence filters 1, 3 and 5 in the \emph{spacetime-average} reading are jointly
unsatisfiable, and the full conjunction of all five filters is the empty class under that reading.
Equivalently: filter~1 is itself the gap test, so ``filter 1 together with computability'' is empty,
and the twenty roads die at filter~1 rather than at the output end.
\end{claimbox}

\src{\code{D06-SPEC.md} \S4.2 (D06-T3 verbatim), \S5.1--\S5.3 (the three steps, with their
acceptance criteria and the memory discipline: one heavy job at a time); \code{SPECIES-EXISTENCE.md}
Thm~2 and Thm~3 (the class-level form and the circularity of lossless quantities).
\LN{} for D06-T3, named class: functionals that are spacetime averages or extrema.}

\subsection{The one route that is still open}
\label{sec:ng-open}

One sufficient condition survives, and it is not a positivity statement. Write
$F(\beta)=\langle\Omega|O|\mathbb G_\beta\rangle$ and
$A(\beta):=\|O^\dagger\Omega_\beta\|\,\|S_W\mathbb G_\beta\|+\|S_W\Omega_\beta\|\,\|O\mathbb G_\beta\|$,
the two terms corresponding to the two terms of $dF/d\beta$. The one-step perturbation expansion gives
$|F'|\le(A/\Delta)|F|$ pointwise, without assuming $F\neq0$, and the first-zero argument then yields

\begin{equation}
(G)\qquad
\int_{\beta_1}^{2.5}\frac{A(\beta)}{\Delta(\beta)}\,d\beta<\infty
\quad\Longrightarrow\quad
f_1^2(2.5)\ \ge\ f_1^2(\beta_1)\exp\Big(-\int_{\beta_1}^{2.5}\frac{A}{\Delta}\,d\beta\Big)\ >\ 0 .
\label{eq:ng-G}
\end{equation}

The logical gain is real: \eqref{eq:ng-G} contains no $F$, so it is not circular, and it consumes only
the \emph{integral} of the reciprocal gap rather than its supremum. It is strictly weaker than a
uniform gap --- a collapse of the form $\Delta(\beta)\ge c/\log(1/|\beta-\beta_c|)$ is allowed, since
$\int1/\Delta=\int\log<\infty$, while $\Delta=c|\beta-\beta_c|^\alpha$ with $\alpha\ge1$ is not --- and
strictly stronger than the target, as a three-level model shows: with
$H=\mathrm{diag}(0,|\beta-\beta_c|^\alpha,1)$, $O=|0\rangle\langle1|+\mathrm{h.c.}$ and
$\Omega=|0\rangle$, one has $F=1>0$ and $F'\equiv0$, yet $(G)$ fails whenever $\alpha\ge1$. So $(G)$ is
sufficient and not necessary, the debt $(G)-\text{target}$ is positive, and whether $(G)$ is even
satisfiable in the physical family --- i.e.\ whether it is vacuous --- is undetermined.
\src{\code{D04-ODE.md} \S1.4 (the one-over-$\Delta$ cost is structural: the diagonal case needs no
$1/\Delta$ but is trivial), \S1.6 (both layers circular without $(G)$), \S3.2 (the implication), \S3.3
(weaker than a uniform gap), \S3.4 (the three-level counterexample; sufficient, not necessary), \S3.5
(verification requires $\Delta>0$ almost everywhere in the window). \LT{} for the implication and for
both comparisons; \LM{} for $(G)$ itself.}

The numbers say how far the condition is from what is provable, and they say it in one line. With the
seed $r_0=1/64$ and $\beta_1\lesssim3.4\times10^{-3}$ (so the same-coordinate Wilson span is
$\Delta\beta_W=2.4848918868$), a floor $F(2.5)\ge\varepsilon$ requires
$C_W\le\log(r_0/\varepsilon)/\Delta\beta_W$: $1.106$ for $\varepsilon=10^{-3}$, $3.886$ for
$10^{-6}$, $9.446$ for $10^{-12}$. The provable constant is
$A_V/\Delta=12L^4/\Delta$, i.e.\ $6.1\times10^3$ at $L=4$, $3.1\times10^4$ at $L=6$ and
$9.8\times10^4$ at $L=8$ --- three to four orders of magnitude too large, and worsening as $L^4$
because $\partial_\beta H=S_W$ is extensive. A multiplicative (Laplace) representation gives
$\log F$ monotone for free with $C\le12V$, hence a floor $e^{-29.96L^4}$: $10^{-3331}$ at $L=4$ and
$10^{-53290}$ at $L=8$, with a premise that is not $L$-uniform. The gap in \eqref{eq:ng-G} is therefore
not an inequality but the $L$-uniformity of $\int A/\Delta$.
\src{\code{D04-ODE.md} \S4.1--\S4.4 (the demand, the provable bound, the multiplicative version,
and the per-$(L,a)$ triviality); \code{Z01-P3ASSEMBLY.md} \S\S4, R4 for $r_0=\tfrac14(2N)^{-2}$ and
$\beta_1$; \code{V05-CPTRANS.md} \S5.4, \S6 for what a discharge would have to be --- a sign carrier in
the $\beta$ direction of identity type, not an inequality --- and for the fact that $\beta_{\mathrm{eff}}$
is defined only from above, so that the interval version of the target is not even well posed yet.}

\medskip
\noindent\textbf{What this chapter does not do.} It does not produce a gap, a bound, or a new route. It
records that the campaign's handles are exhausted in a way that can be stated as a class-level no-go,
that the machine which would replace them is specifiable but not built, and that exactly one
sufficient condition is left unclosed and possibly vacuous. Every claim above is tied to a file and
carries its label; none is marked closed. \textbf{Clay remains $0\%$.}

\let\code\codeprev
